% Курсовая работа — Travel with Ease (XeLaTeX, русский язык)
% 转换自上传文件 «机酒预定网站论文.docx»
\documentclass[14pt,a4paper,oneside]{extreport}
\usepackage{fontspec}
\setmainfont{Liberation Serif}
\setmonofont{Liberation Mono}
\usepackage[russian]{babel}
\usepackage[margin=2cm,left=3cm,right=1.5cm]{geometry}
\usepackage{setspace}
\onehalfspacing
\usepackage{tabularx}
\usepackage{fancyvrb}
\usepackage{hyperref}
\hypersetup{colorlinks=true,linkcolor=black,urlcolor=blue,anchorcolor=black}
\setlength{\parindent}{1.25cm}
\setlength{\parskip}{0pt}
\setlength{\emergencystretch}{3em}
\hyphenpenalty=3000
\tolerance=1200
\begin{document}
\pagestyle{plain}

\begin{titlepage}
\centering
\fontsize{11}{14}\selectfont
МИНИСТЕРСТВО СЕЛЬСКОГО ХОЗЯЙСТВА РОССИЙСКОЙ ФЕДЕРАЦИИ\\
Федеральное государственное Бюджетное образовательное учреждение высшего образования\\
«российский государственный аграрный университет -\\
МСха имени К.А. Тимирязева»\\
(ФГБОУ ВО ргау - МСХА имени К.А. Тимирязева)\\[1.0cm]
{\large Институт экономики и управления\\
Кафедра статистики и кибернетики}\\[0.9cm]
разработка проф. приложений\\[0.9cm]
{\Large\bfseries КУРСОВОЙ ПРОЕКТ}\\[0.4cm]
на тему: Приложение для бронирования гостиниц и поиска\\
авиабилетов с фильтрацией по цене и времени\\[0.9cm]
\begin{flushright}
Выполнила:\\
Студентка 2 курса группы\\
ДЭ 21-25 Хуан Цзини\\
Допущена к защите\\
Руководитель:\\
Калитивин Владимир Анатольевич\\
{\footnotesize ученая степень, ученое звание, ФИО}
\end{flushright}
\vspace{0.8cm}
Оценка \underline{\hspace{6cm}}\\[0.8cm]
\begin{flushright}
Дата \underline{\hspace{3.5cm}}\\
защиты \underline{\hspace{3.5cm}}
\end{flushright}
\vspace{1.0cm}
Москва, 2026
\end{titlepage}
\newpage
\tableofcontents
\newpage
\setcounter{page}{3}

\phantomsection\addcontentsline{toc}{chapter}{Введение}

\chapter*{Введение}

Современное общество характеризуется высоким уровнем мобильности и интенсивным развитием цифровых технологий, которые проникают во все сферы человеческой деятельности. Туризм и путешествия не являются исключением: сегодня планирование поездки практически немыслимо без использования информационных систем, позволяющих в короткие сроки найти актуальную информацию о маршрутах, расписании транспорта, стоимости проживания и погодных условиях.

Цифровизация этой сферы привела к появлению множество онлайн-платформ, которые коренным образом изменили поведение потребителей и саму логику взаимодействия между поставщиками услуг и клиентами.

В условиях постоянного информационного потока, разнообразия клиентских запросов и высокой конкуренции владение навыками проектирования и разработки современных веб-приложений становится не просто желательным, а жизненно необходимым качеством специалиста в области информационных технологий.

Особенно остро эта потребность ощущается в сфере профессионального образования, где студенты готовятся к практической деятельности, предполагающей создание и сопровождение реальных программных продуктов, способных взаимодействовать с пользователями, обрабатывать большие объёмы данных и интегрироваться с внешними сервисами.

Термин «веб-приложение» трактуется в разных источниках неоднозначно, что объясняется его многоплановостью и постоянной эволюцией технических средств. В технической литературе под веб-приложением понимается программное обеспечение, функциональность которого доступна пользователю через сеть Internet при помощи веб-браузера, при этом основная логика обработки данных выполняется на стороне сервера.

наблюдается также подход, согласно которому веб-приложение рассматривается как частный случай клиент-серверной архитектуры, в которой роль клиента выполняет браузер, а роль сервера — веб-сервер с прикладной логикой. В экономическом аспекте веб-приложение выступает как инструмент цифрового взаимодействия, обеспечивающий сокращение транзакционных издержек и расширение доступа пользователей к услугам.

Веб-приложение выполняет множество функций. Прежде всего — информационную, обеспечивающую предоставление актуальных сведений о товарах, услугах и условиях их предоставления. Вторая по значимости функция — транзакционная, позволяющая пользователю совершать операции поиска, сравнения, бронирования и оплаты. Не менее важна коммуникативная функция, проявляющаяся в организации обратной связи, поддержке пользователя и формировании доверия к сервису.

Также выделяются аналитическая и интеграционная функции, обеспечивающие обработку данных и взаимодействие с внешними информационными системами.

Интерес к проблеме проектирования веб-приложений для туристической сферы объясняется тем, что современное общество предъявляет всё более высокие требования к уровню цифровых сервисов. Потребители ожидают от онлайн-платформ не только доступа к информации, но и удобства интерфейса, скорости работы, точности данных и возможности выполнения поиска в любое время и с любого устройства.

В свою очередь, разработчики сталкиваются с необходимостью решать сложные технические задачи: обеспечивать интеграцию с разрозненными источниками данных, корректно обрабатывать ошибки внешних сервисов, оптимизировать скорость ответа и гарантировать защиту персональных данных пользователей.

Актуальность выбранной темы объясняется необходимостью формирования у студентов профессиональных компетенций в области проектирования и разработки веб-приложений, ориентированных на реальные потребности пользователей.

Несмотря на то что вопросы веб-разработки всё чаще включаются в образовательные программы, в практике обучения пока недостаточно практико-ориентированных методик, которые позволяли бы системно связывать теоретические знания с созданием конкретного работающего продукта. Проблема особенно заметна в учреждениях профессионального образования, где выпускники должны быть готовы к самостоятельной разработке и сопровождению информационных систем.

Объектом исследования в данной курсовой работе выступает процесс проектирования и разработки веб-приложения для поиска авиабилетов и бронирования отелей.

Предметом исследования являются методы, архитектурные решения и программные инструменты, обеспечивающие функционирование веб-приложения Travel with Ease, а также принципы интеграции внешних источников данных о рейсах, отелях и погоде.

Цель курсовой работы — теоретически обосновать и практически реализовать веб-приложение для поиска авиабилетов и бронирования отелей, адаптированное для русскоязычных пользователей и функционирующее без необходимости использования платных ключей доступа к внешним сервисам.

Для достижения поставленной цели решаются следующие задачи:

\begin{enumerate}

  \item Изучить научную и техническую литературу по вопросам проектирования веб-приложений и систем онлайн-бронирования, уточнить сущность и структуру основных понятий.

  \item Провести анализ современного состояния рынка цифровых туристических сервисов и сформулировать требования к проектируемому приложению.

  \item Рассмотреть доступные источники данных и программные интерфейсы, выбрать обоснованные решения для реализации функциональных модулей.

  \item Разработать архитектуру и реализовать веб-приложение Travel with Ease, включая модули поиска авиабилетов, бронирования отелей и прогнозирования погоды.

  \item Провести тестирование разработанного приложения и оценить его функциональную пригодность, сформулировать рекомендации по дальнейшему развитию.

\end{enumerate}

Методологическую основу исследования составили системный анализ, методы сравнения и обобщения, изучение технической документации, а также методы проектирования и разработки программного обеспечения с использованием современных инструментов и технологий. В работе применён принцип интеграции открытых данных, позволяющий формировать информационное наполнение приложения без привлечения коммерческих API.

Структура курсовой работы включает введение, две главы, заключение, список литературы и приложение. В первой главе рассматриваются теоретические основы проектирования веб-приложений для поиска авиабилетов и бронирования отелей. Во второй главе описываются процессы проектирования и практической реализации приложения Travel with Ease, а также результаты его тестирования.

\phantomsection\addcontentsline{toc}{chapter}{Глава 1. Теоретические основы проектирования веб-приложений для поиска авиабилетов и бронирования отелей}

\chapter*{Глава 1. Теоретические основы проектирования веб-приложений для поиска авиабилетов и бронирования отелей}

\phantomsection\addcontentsline{toc}{section}{1.1 Сущность и содержание понятий «веб-приложение» и «система онлайн-бронирования»}

\section*{1.1 Сущность и содержание понятий «веб-приложение» и «система онлайн-бронирования»}

Современное развитие информационных технологий привело к появлению нового класса программных продуктов, которые заняли промежуточное положение между традиционными настольными приложениями и статическими веб-страницами. таким продуктом является веб-приложение — программная система, функциональные возможности которой предоставляются пользователю через глобальную сеть Internet с использованием веб-браузера в качестве клиентского интерфейса.

В отличие от обычного веб-сайта, ориентированного преимущественно на представление статической информации, веб-приложение предполагает активное взаимодействие с пользователем, обработку его запросов, изменение состояния данных и формирование персонализированного ответа.

В научной и технической литературе понятие «веб-приложение» определяется с различных позиций. С технической точки зрения веб-приложение представляет собой совокупность программных компонентов, распределённых между клиентской и серверной частями, которые взаимодействуют по протоколу HTTP (HTTPS).

Клиентская часть, или фронтенд, отвечает за отображение интерфейса и обработку действий пользователя; серверная часть, или бэкенд, обеспечивает реализацию бизнес-логики, работу с базой данных и интеграцию с внешними сервисами. С точки зрения пользователя веб-приложение воспринимается как единый интерфейс, не зависящий от установленного программного обеспечения и операционной системы — для работы достаточно наличия браузера и доступа к сети.

Ключевыми характеристиками веб-приложения являются: платформенная независимость, то есть возможность работы на различных устройствах и операционных системах; централизованное хранение данных и логики на сервере, что упрощает обновление и сопровождение; масштабируемость, позволяющая наращивать вычислительные мощности по мере роста нагрузки; а также доступность в любое время и из любого места при наличии соединения с Internet.

Перечисленные качества обусловили широкое распространение веб-приложений в самых разных сферах — от электронной коммерции и банковских услуг до образования и государственного управления.

Особое место среди веб-приложений занимают системы онлайн-бронирования — информационные платформы, предназначенные для поиска, сравнения и резервирования услуг в сфере туризма, транспорта и размещения. Сущность такой системы заключается в предоставлении пользователю возможности самостоятельно, без участия посредника, найти подходящий вариант поездки или проживания, ознакомиться с условиями и ценами, а также оформить предварительную или окончательную заявку.

Современные системы онлайн-бронирования объединяют в себе функции поискового движка, аналитического модуля сравнения предложений, интерфейса оформления заказа и платёжного шлюза.

Структура системы онлайн-бронирования включает несколько функциональных подсистем. Информационная подсистема ответственна за сбор, хранение и актуализацию данных о доступных услугах — рейсах, отелях, тарифах, свободных местах. Поисковая подсистема обеспечивает обработку пользовательских запросов и выборку предложений, соответствующих заданным критериям: дате, маршруту, ценовому диапазону, уровню комфорта.

Подсистема сравнения реализует ранжирование найденных вариантов по заданным параметрам. Транзакционная подсистема поддерживает оформление бронирования, хранение данных заказа и взаимодействие с платёжными системами. Коммуникативная подсистема обеспечивает уведомления пользователя, обратную связь и поддержку.

Важной особенностью систем онлайн-бронирования является их работа с неоднородными данными, поступающими от различных поставщиков. Авиакомпании, операторы систем бронирования, отельные агрегаторы и метеорологические службы используют собственные форматы представления данных, различающиеся по структуре, периодичности обновления и правилам доступа. Это обстоятельство порождает задачу интеграции разрозненных источников — одной из наиболее сложных в проектировании подобных систем.

Эффективное решение этой задачи достигается за счёт использования промежуточного слоя — серверного компонента, который нормализует данные различных источников и предоставляет их клиентскому приложению в едином формате.

Анализ существующих подходов позволяет выделить две основные модели функционирования систем онлайн-бронирования. Первая модель предполагает прямой доступ к собственному массиву данных поставщика: пользователь взаимодействует с платформой, которая хранит актуальные сведения оAvailability номеров или билетов и позволяет совершать бронирование внутри системы.

Вторая модель основана на агрегации — платформа собирает данные из внешних источников и перенаправляет пользователя к поставщику для завершения операции. Каждая из моделей обладает своими достоинствами и ограничениями: прямое бронирование обеспечивает больший контроль над процессом, тогда как агрегация позволяет предложить пользователю более широкий выбор без необходимости заключения договоров с каждым поставщиком.

В контексте данного исследования проектируемое приложение Travel with Ease реализует гибридный подход. Для получения информации о гостиницах и погоде используется прямая интеграция с открытыми источниками данных, результаты которой обрабатываются и отображаются внутри приложения.

Для поиска авиабилетов применяется модель агрегации: приложение формирует структурированный запрос и передаёт пользователя на страницу поиска внешней платформы, где доступны полное расписание и актуальные цены. Такой подход позволяет сочетать преимущества обеих моделей и одновременно избежать необходимости использования платных ключей доступа к коммерческим программным интерфейсам.

Подводя итог рассмотрению понятий, следует отметить, что веб-приложение для поиска авиабилетов и бронирования отелей представляет собой сложную информационную систему, интегрирующую разнообразные источники данных и предоставляющую пользователю единый интерфейс для планирования поездки.

его эффективность определяется не только полнотой функциональных возможностей, но и корректностью обработки данных, скоростью реакции на запросы и удобством взаимодействия, что становится предметом дальнейшего рассмотрения.

\phantomsection\addcontentsline{toc}{section}{1.2 Современное состояние рынка цифровых туристических сервисов}

\section*{1.2 Современное состояние рынка цифровых туристических сервисов}

Рынок цифровых туристических сервисов на сегодняшний день представляет собой одну из наиболее динамично развивающихся сфер цифровой экономики. По данным отраслевых исследований, доля онлайн-каналов в General书中 продаж туристических услуг неуклонно растёт, а количество пользователей, планирующих поездки с помощью цифровых платформ, измеряется сотнями миллионов человек по всему миру.

Эта тенденция обусловлена рядом факторов: распространением мобильных устройств, улучшением качестваINTERNET-соединения, ростом доверия к онлайн-оплате и изменением привычек потребителей, которые всё чаще предпочитают самостоятельно формировать маршрут и выбирать услуги.

Современный рынок можно структурировать по нескольким сегментам. Сегмент поиска и сравнения авиабилетов представлен такими платформами, как Google Flights, Skyscanner, Kayak, Aviasales и другими агрегаторами, которые собирают данные из систем бронирования авиакомпаний и предоставляют пользователю возможность сравнить предложения по цене, времени вылета, количеству пересадок и другим параметрам.

Сегмент бронирования размещения объединяет глобальные платформы (Booking.com, Expedia, Airbnb) и региональные сервисы, предлагающие услуги проживания различных категорий — от хостелов до элитных отелей и апартаментов. Сегмент комплексного планирования поездок включает сервисы, позволяющие одновременно подобрать транспорт, жильё, экскурсии и другие услуги.

Анализ указанных платформ позволяет выявить общие тенденции развития. Во-первых, наблюдается переход от десктопных решений к адаптивным и мобильным интерфейсам, обеспечивающим комфортную работу на смартфонах и планшетах. Во-вторых, платформы стремятся к персонализации предложений за счёт анализа истории поиска и предпочтений пользователя. В-третьих, активно используется искусственный интеллект для прогнозирования цен, рекомендаций маршрутов и автоматизации поддержки.

В-четвёртых, возрастает значение прозрачности данных: пользователи ожидают от сервиса честного отображения цен без скрытых сборов и актуальной информации о наличии мест.

В то же время рынок характеризуется рядом проблемных аспектов, которые создают предпосылки для появления новых решений. Ключевой проблемой является фрагментация данных: информация о рейсах, отелях и погоде распределена между множеством независимых источников, каждый из которых имеет собственный формат и условия доступа.

Многие платформы используют коммерческие программные интерфейсы, требующие регистрации и оплаты, что ограничивает доступ к данным для небольших разработчиков и учебных проектов. Дополнительную сложность создаёт нестабильность данных: расписание рейсов может изменяться, цены — обновляться несколько раз в день, а наличие мест в отелях — зависеть от текущего спроса.

Отдельного внимания заслуживает ситуация на российском рынке туристических сервисов. После изменений в работе международных платформ и ограничения доступа к ряду зарубежных сервисов произошло перераспределение пользовательских потоков в пользу отечественных решений. Платформа Aviasales стала одним из крупнейших агрегаторов авиабилетов в стране, удобный интерфейс поиска и широкий охват направлений.

При этом сохраняется потребность в инструментах, которые дополняли бы функциональность существующих сервисов, обеспечивая, например, отображение реальных объектов размещения на карте или прогноз погоды в месте назначения в увязке с датой поездки.

Конкурентный анализ показывает, что большинство крупных платформ ориентированы на максимизацию продаж и монетизации пользовательского трафика. Это проявляется в активном использовании рекламы, сложных системах тарификации и перенаправлении пользователей на партнёрские сервисы. Для пользователя это означает, что получение полной и объективной информации о поездке сопровождается повышенной когнитивной нагрузкой.

В связи с этим сохраняется ниша для простых, прозрачных приложений, которые не требуют регистрации и предоставляют базовую информацию о маршруте, проживании и погоде в компактной и понятной форме.

Перспективным направлением развития является интеграция открытых данных в пользовательские приложения. Открытые источники, такие как OpenStreetMap для географических объектов и Open-Meteo для метеорологических данных, предоставляют информацию без необходимости использования ключей доступа и без ограничений на коммерческое использование при соблюдении условий лицензий.

Это создаёт технические предпосылки для создания приложений, которые формируют информационное наполнение на основе общедоступных данных, обеспечивая тем самым независимость от коммерческих поставщиков и снижая порог входа для разработчика.

Проведённый анализ рынка позволяет сформулировать несколько важных выводов. Во-первых, спрос на цифровые туристические сервисы продолжает расти, что обусловливает актуальность разработки новых приложений в этой области. Во-вторых, несмотря на разнообразие существующих платформ, сохраняется потребность в простых и прозрачных решениях, ориентированных на русскоязычного пользователя.

В-третьих, использование открытых источников данных является жизнеспособной альтернативой коммерческим интерфейсам, особенно в условиях ограниченного доступа к последним. Перечисленные обстоятельства определили выбор технических решений при проектировании приложения Travel with Ease, которое рассматривается в следующих разделах.

\phantomsection\addcontentsline{toc}{section}{1.3 Обзор источников данных и программных интерфейсов}

\section*{1.3 Обзор источников данных и программных интерфейсов}

Качество функционирования системы онлайн-бронирования в значительной степени определяется полнотой и актуальностью используемых данных. В связи с этим выбор источников информации и программных интерфейсов (API, Application Programming Interface) является одной из ключевых задач проектирования.

В данном разделе рассматриваются основные источники данных, применяемые в приложении Travel with Ease, а также обосновывается их выбор с точки зрения доступности, покрытия и условий использования.

Программный интерфейс приложения представляет собой набор правил и протоколов, позволяющих программным компонентам взаимодействовать друг с другом. В контексте веб-разработки под API обычно понимается интерфейс взаимодействия между клиентским приложением и сервером, а также между сервером приложения и внешними сервисами.

Наиболее распространённым стилем проектирования веб-API является REST (Representational State Transfer), предполагающий использование стандартных HTTP-методов (GET, POST, PUT, DELETE) и Stateless-взаимодействие. Альтернативные подходы включают GraphQL и SOAP, однако для задач интеграции открытых данных чаще всего применяется именно REST благодаря простоте и широкой поддержке.

Для получения информации о рейсах в приложении используется сеть OpenSky Network — глобальная сеть радиолокационных приёмников ADS-B, которая собирает данные о воздушных судах в реальном времени. Через программный интерфейс этой сети предоставляется информация о текущем местоположении самолётов, их высоте, скорости, курсе и позывном. Дополнительно доступна база маршрутов, позволяющая связать позывной рейса с аэропортами вылета и прибытия.

Преимущество OpenSky Network заключается в том, что сервис предоставляет доступ к данным без ключа авторизации для некоммерческого использования, а также в уникальности информации — пользователь может увидеть реальные воздушные суда, находящиеся в воздухе в данный момент. Ограничением является то, что сеть не предоставляет данных о будущих рейсах и ценах билетов, а количество доступных запросов для анонимных пользователей ограничено.

Для получения полного расписания и стоимости авиабилетов применяется подход, основанный на формировании глубинной ссылки (deep link) на страницу поиска платформы Aviasales. Суть подхода состоит в том, что приложение преобразует введённые пользователем данные — города вылета и прибытия, дату и количество пассажиров — в формат, ожидаемый поисковой страницей сервиса, и открывает её в новой вкладке браузера.

Например, строка «MOW1512LED2» кодирует маршрут Москва — Санкт-Петербург, дату 15 декабря и двух пассажиров. Данный подход позволяет предложить пользователю полную информацию о рейсах и ценах без использования платного программного интерфейса и без необходимости регистрации. Код города (IATA-код) определяется либо по встроенному справочнику, либо динамически через сервис подсказок Aviasales.

Информация об отелях формируется на основе данных OpenStreetMap — открытой географической базы данных, создаваемой сообществом пользователей. С помощью сервиса Photon, являющегося поисковым индексом OpenStreetMap, приложение выполняет поиск объектов размещения в заданном городе и получает их реальные названия, адреса и координаты. Это обеспечивает достоверность данных, поскольку информация создаётся и проверяется сообществом, а не генерируется автоматически.

Полученные координаты используются для формирования ссылки на карту и для расчёта расстояния до центра города. Дополнительно приложение формирует ссылку на страницу поиска и цен в системе Booking.com, позволяя пользователю бронирование.

Метеорологические данные предоставляются сервисом Open-Meteo, который на основе современных моделей прогнозирования формирует данные о температуре, осадках, скорости ветра и других параметрах погоды. Сервис поддерживает геокодирование — преобразование названия города в географические координаты — и прогноз погоды на семь дней, что позволяет приложению отображать погоду в городе назначения с подсветкой выбранной пользователем даты.

Open-Meteo не требует ключа доступа и отличается высокой скоростью ответа, что делает его удобным дляintegration в учебные и демонстрационные приложения.

Сравнительный анализ рассмотренных источников приведён в таблице 1.1.

\begin{table}[htbp]\centering

\textbf{Таблица 1.1 — Сравнительная характеристика источников данных}\par\vspace{4pt}

\begin{center}
\small
\begin{tabularx}{\textwidth}{|X|X|X|X|X|}
\hline
Источник & Тип данных & Ключ доступа & Обновление & Назначение в приложении \\
\hline
OpenSky Network & Реальные рейсы в воздухе & Не требуется & Реальное время & Отображение текущих бортов \\
\hline
Aviasales (deep link) & Расписание и цены & Не требуется & По запросу & Поиск всех рейсов на дату \\
\hline
OpenStreetMap / Photon & Отели и объекты размещения & Не требуется & Постоянно & Поиск отелей, карта \\
\hline
Open-Meteo & Погода и геокодирование & Не требуется & Каждые несколько часов & Прогноз погоды \\
\hline
Booking.com & Цены и бронирование & Не требуется (ссылка) & По запросу & Переход к бронированию \\
\hline
\end{tabularx}
\end{center}

\end{table}

Как видно из таблицы, все используемые источники доступны без ключа авторизации, что полностью соответствует требованию независимости от коммерческих сервисов. При этом каждый источник решает собственную задачу, а их интеграция обеспечивает формирование целостного информационного продукта.

Важным аспектом работы с внешними источниками является обработка ошибок и ограничение времени ожидания. Поскольку внешние сервисы могут быть временно недоступны отвечать с задержкой, приложение должно предусматривать механизмы кэширования и таймаутов, а также предоставлять пользователю понятные сообщения об ошибках.

В проектируемом приложении реализован многоуровневый кэш: результаты геокодирования и прогноза погоды хранятся тридцать минут, данные о отелях — десять минут, результаты поиска рейсов — девяносто секунд. Это позволяет снизить нагрузку на внешние сервисы и ускорить повторные запросы.

Подводя итог, следует отметить, что выбранный набор источников данных обеспечивает достаточную полноту информации для решения задач планирования поездки при минимальных затратах на доступ. Использование открытых данных и глубинных ссылок позволяет избежать зависимости от коммерческих ключей и делает приложение воспроизводимым в условиях учебного проекта.

\phantomsection\addcontentsline{toc}{section}{1.4 Принципы проектирования и архитектурные подходы}

\section*{1.4 Принципы проектирования и архитектурные подходы}

Проектирование веб-приложения является многоэтапным процессом, в котором формируются требования, определяется архитектура, разрабатываются компоненты и проверяется соответствие продукта ожиданиям пользователей. В основе этого процесса лежат общие принципы программной инженерии, среди которых выделяются модульность, низкая связность, высокая связность, повторное использование кода и возможность расширения функциональности.

Соблюдение данных принципов особенно важно при создании приложений, интегрирующих несколько внешних источников данных, поскольку позволяет локализовать изменения и упростить сопровождение системы.

Под архитектурой веб-приложения понимается высокоуровневая структура системы, описывающая её основные компоненты, связи между ними и принципы взаимодействия. Наиболее распространённой является трёхзвенная архитектура, включающая клиентский уровень, уровень приложения и уровень данных. Клиентский уровень представлен интерфейсом, работающим в браузере пользователя. Уровень приложения содержит бизнес-логику и обеспечивает взаимодействие с внешними сервисами.

Уровень данных ответственен за хранение информации. В проектируемом приложении функции уровня данных во многом выполняют внешние источники, а локальное хранение ограниченно сводится к сохранению избранных пользователем записей в браузере.

При выборе архитектурного подхода следует учитывать также модель взаимодействия между клиентом и сервером. Традиционная модель предполагает, что сервер формирует готовые HTML-страницы и передаёт их браузеру. Современные приложения чаще используют модель, в которой сервер предоставляет данные в формате JSON, а клиентская часть самостоятельно формирует интерфейс на основе полученных данных.

Такой подход, известный как одностраничное приложение (SPA, Single Page Application), обеспечивает более высокую скорость взаимодействия и позволяет создавать насыщенные интерактивные интерфейсы. В приложении Travel with Ease применён гибридный вариант: сервер отдаёт статическую страницу, а вся динамика интерфейса реализована на языке JavaScript путём обработки ответов серверного API.

Выбор технологического стека является следующей важной задачей проектирования. Для реализации серверной части приложения выбран язык Python и платформа FastAPI, которая ориентирована на создание высокопроизводительных API и обеспечивает автоматическую генерацию документации. К достоинствам FastAPI относятся простота описания маршрутов, поддержка асинхронной обработки запросов и встроенная проверка параметров.

Клиентская часть реализована на основе стандартных технологий HTML, CSS и JavaScript без использования тяжёлых фронтенд-фреймворков, что позволяетKeeper код компактным и понятным, а также упрощает его изучение в рамках профессионального образования.

Значительное внимание в процессе проектирования уделяется адаптивности интерфейса. Поскольку приложение должно быть доступно на устройствах с различным размером экрана, верстка основана на медиазапросах CSS: на мобильных устройствах элементы формы располагаются в один столбец, на экранах бо́льшего размера — в несколько колонок, а навигация адаптируется под формат смартфона.

Такой подход обеспечивает единообразие пользовательского опыта независимо от используемого устройства.

Принцип доступности и удобства использования реализуется также через отказ от обязательной регистрации пользователя. Приложение не требует создания учётной записи, а сохранение избранных элементов осуществляется локально в хранилище браузера. Это снижает порог входа и соответствует задаче создания прозрачного сервиса, ориентированного на быстрое получение информации.

В то же время для совершения бронирования пользователь направляется на страницу внешнего проверенного сервиса, что позволяет избежать хранения конфиденциальных платёжных данных внутри приложения.

Особое место в архитектуре занимает слой интеграции с внешними источниками. Поскольку каждый источник имеет собственный формат данных, серверное приложение выполняет функции адаптера: оно отправляет запросы, получает ответы, нормализует структуру данных и формирует единый ответ для клиента. Такой подход позволяет скрыть неоднородность источников от клиентской части и упрощает замену или добавление новых источников в будущем.

Кроме того, на этом уровне реализованы кэширование и обработка ошибок, о которых говорилось в предыдущем разделе.

Проектирование данных приложения ориентировано на минимализм. Основными сущностями являются город (с названием и кодом), рейс (с позывным, авиакомпанией, аэропортами и parametрами полёта), отель (с названием, адресом и координатами) и погодный прогноз (с датой и meteorологическими показателями). Каждая сущность представлена в виде структуры данных, передаваемой клиенту в формате JSON.

Отсутствие собственной реляционной базы данных объясняется тем, что основная информация поступает из внешних источников в режиме реального времени, а постоянному хранению подлежит лишь небольшой объём пользовательских предпочтений.

Оценка качества архитектуры осуществляется по ряду критериев: производительность, надёжность, масштабируемость, удобство сопровождения и безопасность.

Производительность обеспечивается за счёт кэширования и параллельной обработки запросов; надёжность — за счёт обработки ошибок и таймаутов; масштабируемость — за счёт модульной структуры; удобство сопровождения — за счёт разделения ответственности между компонентами; безопасность — за счёт отсутствия хранения конфиденциальных данных и использования защищённого протокола при обмене с внешними сервисами.

Таким образом, теоретический анализ, проведённый в первой главе, позволяет сформулировать следующие выводы. Веб-приложение для поиска авиабилетов и бронирования отелей является актуальным типом информационной системы, востребованным на современном рынке цифровых услуг. Существующие платформы, несмотря на разнообразие функциональных возможностей, оставляют нишу для простых и прозрачных решений, ориентированных на русскоязычного пользователя.

Использование открытых источников данных и глубинных ссылок позволяет реализовать такое решение без привлечения коммерческих ключей. Выбранные архитектурные подходы и технологический стек обеспечивают возможность praktической реализации приложения, которая описывается во второй главе.

\phantomsection\addcontentsline{toc}{chapter}{Глава 2. Проектирование и разработка веб-приложения Travel with Ease}

\chapter*{Глава 2. Проектирование и разработка веб-приложения Travel with Ease}

\phantomsection\addcontentsline{toc}{section}{2.1 Общая архитектура и структура приложения}

\section*{2.1 Общая архитектура и структура приложения}

Практическая реализация приложения Travel with Ease осуществлена на основе трёхзвенной архитектуры, описанной в предыдущей главе. В качестве серверной платформы используется язык Python версии 3.11 и веб-фреймворк FastAPI, который обеспечивает обработку HTTP-запросов, проверку параметров и формирование ответов в формате JSON.

Клиентская часть представляет собой одностраничный интерфейс, построенный на технологиях HTML5, CSS3 и JavaScript, что позволяет отказаться от дополнительных зависимостей и обеспечить быструю загрузку страницы. Взаимодействие между клиентом и сервером осуществляется через программный интерфейс, предоставляющий несколько конечных точек.

Структура проекта организована следующим образом. Корневой каталог содержит файл приложения, в котором определены маршруты и бизнес-логика, а также каталог статических файлов с клиентским интерфейсом. Серверное приложение при запуске монтирует каталог статических файлов и отдаёт главную страницу по корневому адресу. Такая организация позволяет разворачивать приложение как единое целое и упрощает его перенос на хостинг-платформу.

Основные конечные точки программного интерфейса приложения перечислены в таблице 2.1.

\begin{table}[htbp]\centering

\textbf{Таблица 2.1 — Основные конечные точки API приложения}\par\vspace{4pt}

\begin{center}
\small
\begin{tabularx}{\textwidth}{|X|X|X|}
\hline
Маршрут & Метод & Назначение \\
\hline
/ & GET & Отдача главной страницы интерфейса \\
\hline
/api/flights & GET & Поиск реальных рейсов в воздухе (OpenSky) \\
\hline
/api/flight\_link & GET & Формирование ссылки на поиск всех рейсов (Aviasales) \\
\hline
/api/hotels & GET & Поиск отелей в городе (OpenStreetMap) \\
\hline
/api/weather & GET & Получение прогноза погоды \\
\hline
\end{tabularx}
\end{center}

\end{table}

Каждый маршрут реализует собственную функцию-обработчик, которая получает параметры запроса, обращается к соответствующему внешнему источнику, обрабатывает результат и возвращает структурированный ответ. Обработка параметров включает проверку обязательности полей, ограничение длины строк и диапазона числовых значений.

В случае некорректных данных сервер возвращает сообщение об ошибке с соответствующим кодом состояния HTTP, что позволяет клиентской части корректно реагировать на ситуацию.

Клиентский интерфейс разделён на три функциональные области, соответствующие основным задачам приложения: поиск авиабилетов, поиск отелей и избранное. Переключение между областями осуществляется без перезагрузки страницы путём изменения класса активного элемента. Такой подход обеспечивает мгновенную реакцию интерфейса и создаёт ощущение целостного приложения.

На мобильных устройствах для переключения используется нижняя панель навигации, а на больших экранах — горизонтальные вкладки.

Организация обмена данными между клиентом и сервером построена на асинхронных запросах с использованием конструкции fetch и механизма Promise. Это позволяет интерфейсу оставаться отзывчивым во время ожидания ответа от внешних сервисов. На время выполнения запроса кнопка поиска блокируется и отображается индикатор загрузки, что информирует пользователя о ходе операции. После получения ответа данные форматируются и отображаются в области результатов.

Важным элементом архитектуры является слой кэширования. Как отмечалось в разделе 1.3, для хранения результатов запросов используется словарь с временными метками. Функция кэширования проверяет наличие актуальной записи и, если она найдена, возвращает её без обращения к внешнему источнику. Это особенно важно для сервисов с ограничением частоты запросов, к которым относится OpenSky Network.

Реализация кэша не требует внешней базы данных и размещается в памяти процесса приложения.

Обработка ошибок реализована на двух уровнях. На серверном уровне обрабатываются исключения, возникающие при обращении к внешним сервисам, и формируются понятные сообщения. На клиентском уровне используется конструкция try-catch, которая перехватывает ошибки сети и отображает диагностическое сообщение.

Отдельно предусмотрена ситуация, когда страница открыта как локальный файл: в этом случае запросы к серверному интерфейсу невозможны, и пользователь получает соответствующее предупреждение.

Локализация интерфейса выполнена на русском языке. Все тексты элементов управления, сообщения и форматы дат соответствуют русскоязычной норме: для форматирования дат используется метод toLocaleDateString с локалью «ru-RU», а числовые значения форматируются с использованием разделителя разрядов. Это обеспечивает естественное восприятие интерфейса целевой аудиторией.

Подводя итог описанию архитектуры, можно отметить, что приложение представляет собой компактную систему с чётким разделением ответственности между компонентами. Использование современных инструментов и отказ от избыточных зависимостей позволили создать решение, которое легко разворачивается, быстро работает и может быть адаптировано для учебных целей. Детальное описание отдельных модулей приводится в следующих разделах.

\phantomsection\addcontentsline{toc}{section}{2.2 Реализация модуля поиска авиабилетов}

\section*{2.2 Реализация модуля поиска авиабилетов}

Модуль поиска авиабилетов является одной из центральных функций приложения. В процессе проектирования была поставлена задача предоставить пользователю возможность получения полной информации о рейсах на выбранную дату — включая расписание, время вылета и прилёта, авиакомпанию и стоимость — без использования платного программного интерфейса.

Как было показано в разделе 1.3, такая возможность реализуется посредством формирования глубинной ссылки на страницу поиска платформы Aviasales.

Интерфейс модуля содержит поля ввода города вылета и города прибытия, поле выбора даты, элемент выбора количества пассажиров, а также кнопку поиска. Для городов реализован механизм подсказок, упрощающий ввод. Дополнительно предусмотрена кнопка быстрой смены направления и набор ссылок на популярные маршруты, что сокращает количество действий пользователя.

При нажатии кнопки поиска клиентская часть выполняет запрос к маршруту /api/flight\_link, передавая параметры from\_city, to\_city, date и adults. Серверная часть выполняет несколько последовательных операций. Во-первых, производится преобразование названий городов в IATA-коды. Для этого используется встроенный справочник, содержащий коды основных российских и международных городов.

Если город отсутствует в справочнике, выполняется обращение к сервису подсказок Aviasales, который возвращает код города по его названию. Результат поиска кэшируется на двадцать четыре часа, что исключает повторные обращения при аналогичных запросах.

Во-вторых, сервер проверяет корректность даты и преобразует её в формат, ожидаемый поисковой страницей. Строка даты вида «2026-12-15» преобразуется в комбинацию «1512», где первые две цифры обозначают день, а последние две — месяц. Затем формируется итоговый адрес вида \url{https://www.aviasales.ru/search/MOW1512LED2,} где после кода города вылета следует дата, затем код города прибытия и количество пассажиров.

В-третьих, сервер дополнительно запрашивает прогноз погоды в городе прибытия на выбранную дату, чтобы отобразить его в карточке результата. Если получить прогноз не удаётся, формируется без погодной информации, что не нарушает основную функциональность. Итоговый ответ содержит названия городов, их коды, дату, количество пассажиров, сформированную ссылку и данные погоды.

Клиентская часть получает ответ и формирует карточку результата, которая включает маршрут, дату поездки, коды аэропортов, блок прогноза погоды с подсветкой выбранной даты и две кнопки: переход к поиску и сохранение запроса в избранное. Переход осуществляется в новой вкладке браузера, чтобы пользователь мог вернуться к приложению. Дляundi случаев блокировки всплывающих окон ссылка доступна как обычная гиперссылка.

Следует отдельно обсудитьQuestion точности поиска. На предыдущих этапах разработки приложение использовало данные сети OpenSky для отображения реальных воздушных судов, находящихся в воздухе. Однако этот подход имел существующий недостаток: результаты зависели от текущего времени суток и не позволяли получить расписание на будущую дату.

Кроме того, при попытке расширить выдачу за счёт частичного совпадения ( когда совпадал только город вылета или только город прибытия) в результаты попадали рейсы в третьи города, что вводило пользователя в заблуждение. Поэтому в итоговой версии для поиска всех рейсов применяется модель глубинной ссылки, а отображение реальных бортов сохранено как дополнительная функция на основе OpenSky с гарантией полного совпадения направления.

Сравнение двух подходов к поиску авиабилетов приведено в таблице 2.2.

\begin{table}[htbp]\centering

\textbf{Таблица 2.2 — Сравнение подходов к поиску авиабилетов}\par\vspace{4pt}

\begin{center}
\small
\begin{tabularx}{\textwidth}{|X|X|X|}
\hline
Характеристика & OpenSky (реальное время) & Глубинная ссылка Aviasales \\
\hline
Покрытие по времени & Только текущий момент & Любая выбранная дата \\
\hline
Стоимость & Бесплатно & Бесплатно \\
\hline
Ключ доступа & Не требуется & Не требуется \\
\hline
Цены билетов & Отсутствуют & Присутствуют \\
\hline
Возможность бронирования & Нет & Да \\
\hline
Зависимость от времени суток & Высокая & Отсутствует \\
\hline
\end{tabularx}
\end{center}

\end{table}

Как следует из таблицы, глубинная ссылка превосходит подход OpenSky по большинству показателей, актуальных для конечного пользователя. Именно поэтому он выбран в качестве основного механизма поиска.

Особое внимание уделено корректности кодирования городов. В справочник включены кодыMosCOW — MOW, Санкт-Петербург — LED, Сочи — AER, Екатеринбург — SVX, Казань — KZN и другие. Для международных направлений указаны коды Дубай — DXB, Стамбул — IST, Бангкок — BKK и прочие. Использование кодов городов, а не конкретных аэропортов, позволяет охватить все аэропортыcorresponding города — например, код MOW включает Шереметьево, Домодедово и Внуково.

Функция сохранения запроса в избранное реализована с использованием локального хранилища браузера localStorage. Каждая запись содержит тип, временную метку и объект с параметрами поиска. При отображении списка избранного для каждого запроса формируется ссылка, позволяющая повторить поиск в один клик. Ограничение на fiftyть записей предотвращает неконтролируемый рост хранилища.

Таким образом, модуль поиска авиабилетов обеспечивает решение основной задачи — получение полной информации о рейсах и ценах на выбранную дату без регистрации и ключей доступа. Реализация отличается простотой, надёжностью и удобством использования.

\phantomsection\addcontentsline{toc}{section}{2.3 Реализация модуля поиска отелей и погодного сервиса}

\section*{2.3 Реализация модуля поиска отелей и погодного сервиса}

Модуль поиска отелей предназначен для предоставления пользователю информации о местах размещения в выбранном городе. В основе модуля лежит использование данных OpenStreetMap посредством поискового сервиса Photon. Серверный обработчик маршрута /api/hotels получает название города, выполняет геокодирование для определения координат центра и затем направляет запрос к сервису Photon с фильтром по типу объекта «tourism=hotel».

Результатом является перечень объектов размещения с реальными названиями, адресами и координатами.

После получения списка сервер выполняет его обработку. Для каждого объекта извлекаются название, улица, номер дома, координаты и тип размещения. На основе координат рассчитывается расстояние до центра города по формуле гаверсинуса, после чего список сортируется по возрастанию расстояния. Это позволяет пользователю сначала увидеть объекты, расположенные ближе к центру.

Для каждого отеля формируется ссылка на карту OpenStreetMap с точкой в координатах объекта и ссылка на страницу поиска цен и бронирования в системе Booking.com с учётом дат заезда и выезда.

Интерфейс\_MODULE позволяет пользователю указать город, дату заезда и дату выезда. На основе дат автоматически вычисляется количество ночей, что отображается в описании каждого отеля и передаётся в ссылку бронирования. Проверка корректности дат исключает ситуацию, когда дата выезда предшествует дате заезда. Результаты отображаются в виде карточек с названием, типом размещения, адресом, расстоянием до центра и набором кнопок действий.

Погодный сервис реализован на основе Open-Meteo и используется в обоих модулях — при поиске авиабилетов и при поиске отелей. Серверная функция геокодирования получает координаты города, а функция прогноза запрашивает данные о максимальной и минимальной температуре, вероятности осадков и коде погодного явления на семь дней вперёд. Код явления преобразуется в текстовое описание и эмодзи, что делает интерфейс наглядным.

Текущая погода отображается в заголовке блока вместе с датой, а выбранная пользователем дата подсвечивается среди дней прогноза.

Стоит отметить, что прогноз доступен только на семь дней вперёд. Если пользователь выбирает дату, выходящую за пределы прогноза, интерфейс отображает соответствующее пояснение. Это решение принято для того, чтобы не вводить пользователя в заблуждение относительно доступности данных. Значение выбранной даты сохраняется в карточке результата и используется как элемент планирования поездки.

В таблице 2.3 приведены основные поля данных отеля, используемые в приложении.

\begin{table}[htbp]\centering

\textbf{Таблица 2.3 — Структура данных об объекте размещения}\par\vspace{4pt}

\begin{center}
\small
\begin{tabularx}{\textwidth}{|X|X|X|}
\hline
Поле & Описание & Пример \\
\hline
name & Название объекта & Сочи Марриотт Красная Поляна \\
\hline
street & Улица & улица Заречная \\
\hline
housenumber & Номер дома & 1 \\
\hline
kind & Тип размещения & hotel, guest\_house, hostel \\
\hline
lat, lon & Географические координаты & 43.60, 39.73 \\
\hline
dist\_km & Расстояние до центра, км & 2.4 \\
\hline
map\_url & Ссылка на карту & openstreetmap.org/... \\
\hline
\end{tabularx}
\end{center}

\end{table}

Обработка ошибок в модуле построена по тому же принципу, что и в модуле авиабилетов. Если город не найден или внешний сервис недоступен, пользователь получает понятное сообщение. Благодаря кэшированию повторные запросы к одному и тому же городу выполняются мгновенно, что положительно сказывается на опыте взаимодействия.

Результатом работы модуля является список реальных объектов размещения с достоверными адресами и координатами, дополненный прогнозом погоды и возможностью перехода к бронированию. Это полностью соответствует требованиям, сформулированным в первой главе, и обеспечивает практическую ценность приложения.

\phantomsection\addcontentsline{toc}{section}{2.4 Тестирование и оценка эффективности приложения}

\section*{2.4 Тестирование и оценка эффективности приложения}

Завершающим этапом разработки является тестирование приложения и оценка его соответствия предъявленным требованиям. Тестирование проводилось в несколько направлений: проверка синтаксиса исходного кода, функциональное тестирование каждого модуля, проверка корректности обработки пользовательского ввода и оценка адаптивности интерфейса.

На первом этапе была выполнена проверка синтаксиса серверного кода и клиентского сценария. Для серверной части использовалась стандартная компиляция Python, для клиентской — синтаксическая проверкаJavaScript с помощью среды Node.js. Ошибки синтаксиса не обнаружены, что подтверждает корректность структуры исходного кода.

На втором этапе проводилось функциональное тестирование программного интерфейса. С помощью инструмента командной строкиcurl и языка Python выполнялись запросы к каждому маршруту с последующим анализом ответов. Проверялись следующие сценарии: поиск отелей по названию города, получение прогноза погоды, формирование ссылки на поиск авиабилетов для известного и неизвестного города, а также обработка некорректных параметров. Результаты тестирования представлены в таблице 2.4.

\begin{table}[htbp]\centering

\textbf{Таблица 2.4 — Результаты функционального тестирования}\par\vspace{4pt}

\begin{center}
\small
\begin{tabularx}{\textwidth}{|X|X|X|}
\hline
Тестовый сценарий & Ожидаемый результат & Фактический результат \\
\hline
Поиск отелей «Сочи» & Список реальных отелей & Успешно, найдено несколько объектов \\
\hline
Прогноз погоды «Сочи» & Данные на 7 дней & Успешно, получены данные \\
\hline
Ссылка Москва → Санкт-Петербург & Коды MOW → LED & Успешно, сформирована ссылка \\
\hline
Ссылка для неизвестного города & Определение кода через подсказки & Успешно, код определён \\
\hline
Одинаковые города вылета и прибытия & Сообщение об ошибке & Успешно, ошибка возвращена \\
\hline
Запуск через локальный файл & Предупреждение о недоступности & Успешно, предупреждение отображено \\
\hline
\end{tabularx}
\end{center}

\end{table}

На третьем этапе использовался инструмент автоматизированного браузерного тестирования, который позволил проверить работу интерфейса в реальных условиях. Были выполнены следующие действия: открытие страницы, выбор городов и дат, нажатие кнопок поиска, проверка появления карточек результатов, сохранение в избранное и переключение между разделами. В ходе тестирования фиксировались ошибки выполнения сценария JavaScript.

Ошибок не обнаружено, все элементы интерфейса работали корректно.

Отдельно оценивалась скорость работы приложения. При первом обращении к внешним сервисам время ответа составляло от двух до пяти секунд, что обусловлено необходимостью получения данных из нескольких источников. При повторных запросах благодаря кэшированию время ответа сокращалось до долей секунды. Это подтверждает эффективность реализованного механизма кэширования.

Оценка адаптивности интерфейса проводилась путём изменения размера окна браузера. На узком экране поля формы располагались в один столбец, элементы управления увеличивались в размере для удобства нажатия, а навигация отображалась в нижней части экрана. На широком экране интерфейс принимал вид с центральным содержимым и горизонтальным меню. Проблемы наложения или обрезки элементов не выявлены.

По результатам тестирования были сформулированы следующие выводы. Приложение соответствует функциональным требованиям, предъявленным на этапе проектирования. Все основные модули работают корректно, обработка ошибок обеспечивает устойчивость к недоступности внешних сервисов, а кэширование обеспечивает приемлемую скорость работы. Интерфейс адаптирован для использования на мобильных и десктопных устройствах.

В то же время выявлен ряд направлений для дальнейшего развития: расширение справочника городов, добавление поддержки обратного направления поиска, интеграция дополнительных источников данных и реализация собственной базы данных для хранения истории поиска.

Комплексная оценка приложения по критериям, сформулированным в разделе 1.4, приведена в таблице 2.5.

\begin{table}[htbp]\centering

\textbf{Таблица 2.5 — Оценка приложения по критериям качества}\par\vspace{4pt}

\begin{center}
\small
\begin{tabularx}{\textwidth}{|X|X|X|}
\hline
Критерий & Оценка & Комментарий \\
\hline
Функциональность & Высокая & Все заявленные функции реализованы \\
\hline
Производительность & Высокая & Кэширование обеспечивает быстрый ответ \\
\hline
Надёжность & Хорошая & Обработка ошибок внешних сервисов \\
\hline
Адаптивность & Высокая & Корректная работа на разных устройствах \\
\hline
Удобство использования & Высокая & Отсутствие регистрации, понятный интерфейс \\
\hline
Сопровождаемость & Хорошая & Модульная структура, минимальные зависимости \\
\hline
\end{tabularx}
\end{center}

\end{table}

Таким образом, результаты тестирования подтверждают достижение цели курсовой работы и возможность практического использования разработанного приложения.

\phantomsection\addcontentsline{toc}{section}{2.5 Проектирование структур данных и алгоритмов преобразования информации}

\section*{2.5 Проектирование структур данных и алгоритмов преобразования информации}

Проектирование структур данных является ключевым этапом создания любого информационного приложения, поскольку от выбранной модели представления данных зависят скорость обработки запросов, удобство расширения функционала и надёжность взаимодействия между отдельными модулями. В разработанном приложении Travel with Ease данные получаются из нескольких внешних источников, каждый из которых использует собственный формат описания объектов.

Поэтому была поставлена задача сформировать единые внутренние структуры данных, которые скрывали бы различия между источниками и предоставляли клиентскому интерфейсу однородный набор полей.

При проектировании внутренних структур был применён принцип минимальной достаточности: в состав каждой модели включались только те поля, которые действительно используются в интерфейсе или в алгоритмах обработки. Это позволило снизить объём передаемых данных и упростить сериализацию ответов серверного интерфейса. Основные модели приложения включают описание рейса, описание объекта размещения и описание погодных условий.

Каждая модель реализована в виде словаря (объекта JSON), что обеспечивает прямую совместимость с форматом обмена данными REST-интерфейса.

Модель рейса содержит поля, характеризующие воздушное судно и его полёт: позывной, наименование авиакомпании, dep\_name и arr\_name (наименования аэропортов вылета и прибытия), dep\_icao и arr\_icao (коды ICAO), а также координаты, высоту, скорость и степень завершённости маршрута. Для режима поиска по дате в модель добавляются поля с кодом города, датой поездки и сформированной ссылкой на внешнюю платформу.

Такое разделение позволяет использовать одну и ту же модель в обоих режимах поиска без дублирования кода.

Модель объекта размещения формируется на основе данных OpenStreetMap и включает наименование, тип размещения, улицу, номер дома, географические координаты и расстояние до центра города. Расстояние вычисляется по формуле гаверсинуса, которая устойчиво работает на больших расстояниях и не требует привлечения дополнительных библиотек.

Модель погоды содержит текущую температуру, текстовое описание состояния погоды, а также массив прогнозов по дням с максимальной и минимальной температурой и вероятностью осадков.

Отдельным объектом проектирования стал алгоритм преобразования наименования города в код IATA. Алгоритм работает в два этапа. На первом этапе выполняется поиск в локальной таблице соответствий, содержащей наиболее часто используемые российские и международные направления. Если город найден, код возвращается немедленно без обращения к сети.

На втором этапе при отсутствии соответствия выполняется запрос к сервису подсказок Aviasales, из ответа которого извлекается код первого объекта типа «город». Такой гибридный подход обеспечивает высокую скорость для популярных направлений и одновременно поддерживает произвольные города без необходимости вручную обновлять справочник.

Алгоритм формирования глубинной ссылки также реализован в несколько этапов. Исходная дата в формате ГГГГ-ММ-ДД преобразуется в компактный формат ДДММ, после чего выполняется объединение кодов городов и числа пассажиров. Полученная строка включается в шаблон адреса поисковой страницы. Например, дата 2026-12-15, города Москва и Санкт-Петербург и два пассажира дают ссылку вида MOW1512LED2.

Достоинством подхода является полная детерминированность: одинаковые входные данные всегда дают одинаковый адрес, что упрощает тестирование и кэширование результатов.

Важным элементом проектирования стала система кэширования. Поскольку внешние сервисы имеют ограничения на частоту запросов и могут отвечать с задержкой, были введены три уровня кэша: результаты поиска рейсов хранятся девяносто секунд, результаты поиска отелей — десять минут, а данные геокодирования и погоды — от тридцати минут до двадцати четырёх часов.

Повторный запрос к тому же маршруту в пределах срока действия кэша обрабатывается практически мгновенно, что подтверждено результатами тестирования.

Обработка ошибок также была включена в алгоритмы преобразования. Если внешний сервис недоступен или возвращает некорректный ответ, серверный модуль формирует диагностическое сообщение, которое отображается в интерфейсе вместо пустой страницы. Аналогично обрабатываются ситуации совпадения городов вылета и прибытия, пустого ввода и открытия страницы в обход сервера. Таким образом, пользователь всегда получает понятное объяснение произошедшей ошибки.

\begin{table}[htbp]\centering

\textbf{Таблица 2.6 — Основные внутренние структуры данных приложения}\par\vspace{4pt}

\begin{center}
\small
\begin{tabularx}{\textwidth}{|X|X|X|}
\hline
Модель & Ключевые поля & Источник данных \\
\hline
Рейс в воздухе & позывной, авиакомпания, dep\_icao, arr\_icao, lat, lon, alt\_m & OpenSky Network \\
\hline
Рейс на дату & from\_code, to\_code, date, adults, url & Aviasales (глубинная ссылка) \\
\hline
Объект размещения & name, street, lat, lon, dist\_km, kind & OpenStreetMap / Photon \\
\hline
Погода & current, daily, rain\_prob, emoji & Open-Meteo \\
\hline
\end{tabularx}
\end{center}

\end{table}

Подводя итог рассмотрению структур данных и алгоритмов, следует отметить, что применение единых внутренних моделей, гибридного алгоритма определения кодов городов и многоуровневого кэширования позволило обеспечить согласованную работу всех модулей приложения. Принятые решения создают основу для дальнейшего расширения функциональности, например для добавления собственной базы данных или поддержки нескольких языков интерфейса.

\phantomsection\addcontentsline{toc}{section}{2.6 Проектирование адаптивного пользовательского интерфейса и реализация пользовательских сценариев}

\section*{2.6 Проектирование адаптивного пользовательского интерфейса и реализация пользовательских сценариев}

Пользовательский интерфейс является той частью приложения, с которой взаимодействует конечный пользователь, поэтому его проектирование получило особое внимание. В ходе разработки была поставлена задача создать простой и понятный интерфейс, который был бы одинаково удобен на мобильном устройстве и на персональном компьютере.

Для достижения этой цели была применена методология mobile-first, при которой сначала проектируется компактная версия для небольшого экрана, а затем с помощью медиазапросов добавляются стили для больших размеров экрана.

Структура интерфейса организована на трёх вкладках: «Авиабилеты», «Отели» и «Избранное». Такое решение позволяет избежать перегрузки экрана и дать пользователю возможность быстро переключаться между задачами. На мобильных устройствах вкладки размещаются в нижней части экрана в виде панели навигации, что соответствует привычным паттернам взаимодействия со смартфоном.

На компьютере вкладки располагаются в верхней части страницы, а содержимое размещается в центре экрана с ограничением максимальной ширины.

Форма поиска авиабилетов включает поля города вылета, города прибытия, даты поездки и количества пассажиров. Для ускорения ввода реализованы кнопки популярных направлений, такие как Москва — Сочи и Москва — Санкт-Петербург. Кнопка обмена городами позволяет одним нажатием изменить направление поиска.

После нажатия кнопки поиска серверный модуль формирует ссылку, а в интерфейсе отображается карточка результата с маршрутом, датой, кодом городов и прогнозом погоды в месте назначения. Крупная кнопка перехода открывает страницу поисковой системы в новом окне.

Форма поиска отелей включает город назначения и даты заезда и выезда. Количество ночей вычисляется автоматически как разность между датами. Результат отображается в виде списка карточек, каждая из которых содержит наименование объекта, тип размещения, адрес, расстояние до центра и ссылки на карту и на страницу бронирования. Ссылка на бронирование формируется с учётом выбранных дат, что позволяет пользователю сразу перейти к подходящим предложениям на даты его поездки.

Особое внимание было уделено визуальным состояниям интерфейса. Во время выполнения запроса кнопка поиска блокируется и в ней отображается индикатор загрузки, что предотвращает повторные запросы и информирует пользователя о ходе операции. При возникновении ошибки вместо результата показывается диагностическое сообщение с объяснением причины. Если поиск выполнен успешно, но данные отсутствуют, отображается специальное пустое состояние с рекомендацией изменить дату или город.

Функция избранного реализована с использованием локального хранилища браузера. Сохранённые запросы и объекты отображаются в отдельной вкладке и содержат те же поля, что и карточки результатов, включая дату и ссылку. Это позволяет пользователю вернуться к ранее найденному варианту без повторного поиска. Важно, что данные хранятся только на устройстве пользователя и не передаются на сервер, что снижает нагрузку на инфраструктуру и упрощает соблюдение конфиденциальности.

Адаптивность интерфейса обеспечивается за счёт относительных единиц измерения, гибкой сетки и медиазапросов. На узких экранах элементы формы располагаются в один столбец, а карточки результатов занимают всю доступную ширину. На широких экранах форма отображается в виде сетки из нескольких колонок, а результаты центрируются по ширине страницы. Шрифты и отступы подобраны так, чтобы сохранить читаемость на всех размерах экрана.

\begin{table}[htbp]\centering

\textbf{Таблица 2.7 — Соответствие пользовательских сценариев элементам интерфейса}\par\vspace{4pt}

\begin{center}
\small
\begin{tabularx}{\textwidth}{|X|X|X|}
\hline
Сценарий & Элемент интерфейса & Результат \\
\hline
Поиск рейсов на дату & Вкладка «Авиабилеты», форма поиска & Карточка со ссылкой и погодой \\
\hline
Поиск отелей & Вкладка «Отели», даты заезда и выезда & Список карточек с ценами \\
\hline
Сохранение варианта & Кнопка «В избранное» & Запись в локальном хранилище \\
\hline
Просмотр избранного & Вкладка «Избранное» & Список сохранённых объектов \\
\hline
\end{tabularx}
\end{center}

\end{table}

Таким образом, спроектированный интерфейс полностью покрывает пользовательские сценарии, определённые на этапе анализа требований. Использование адаптивной вёрстки, понятных состояний и локального хранилища позволяет считать приложение удобным для повседневного использования на различных устройствах.

\phantomsection\addcontentsline{toc}{chapter}{Заключение}

\chapter*{Заключение}

Заключение курсовой работы формирует целостный итог проделанной работы и позволяет показать, как теоретические положения, проектные решения и практическая реализация объединяются в единый процесс создания веб-приложения для поиска авиабилетов и бронирования отелей.

Развитие навыков проектирования информационных систем, работы с внешними источниками данных и создания удобных интерфейсов становится для современных обучающихся не менее важным, чем фундаментальная теоретическая подготовка. В условиях, когда рынок труда требует от молодых специалистов не только владения технологиями, но и способности создавать востребованные программные продукты, компетенции в области веб-разработки приобретают статус базовых.

В ходе курсовой работы были последовательно раскрыты ключевые подходы к пониманию сущности веб-приложения и системы онлайн-бронирования. Показано, что веб-приложение представляет собой распределённую информационную систему, в которой клиентская часть обеспечивает взаимодействие с пользователем, а серверная — реализацию бизнес-логики и интеграцию с внешними источниками.

Теоретический анализ продемонстрировал, что современный рынок цифровых туристических сервисов характеризуется высоким уровнем конкуренции и фрагментацией данных, что создаёт предпосылки для появления простых и прозрачных приложений, ориентированных на русскоязычного пользователя.

Итогом теоретического осмысления стал конкретный программный продукт — веб-приложение Travel with Ease. Его ценность заключается не в новизне используемых технологий самих по себе, а в продуманной интеграции открытых источников данных и глубинных ссылок, позволяющей получить полную информацию о поездке без использования платных ключей доступа.

Приложение объединяет три функциональных модуля: поиск авиабилетов с переходом к полному расписанию и ценам, поиск реальных объектов размещения на основе данных OpenStreetMap и прогноз погоды в месте назначения.

Реализация приложения выполнена на основе трёхзвенной архитектуры с использованием платформы FastAPI и современных клиентских технологий. В ходе разработки были решены задачи геокодирования городов, преобразования дат в формат глубинной ссылки, нормализации данных внешних источников, кэширования запросов и обработки ошибок.

Особое внимание уделено адаптивности интерфейса, что позволяет использовать приложение на мобильных устройствах и персональных компьютерах без потери функциональности.

Практическая значимость работы подтверждена результатами тестирования. Проверка синтаксиса, функциональное тестирование маршрутов серверного интерфейса и автоматизированное браузерное тестирование показали корректность работы всех модулей. Ошибки выполнения сценария JavaScript не обнаружены. Время повторных запросов благодаря кэшированию сокращается до долей секунды.

Приложение успешно обрабатывает некорректные пользовательские вводы и недоступность внешних сервисов, предоставляя понятные диагностические сообщения.

Вместе с тем в ходе работы выявлен ряд ограничений, которые определяют направления дальнейшего развития. Прежде всего, прогноз погоды доступен только на семь дней вперёд, что связано с возможностями используемого сервиса. Поиск авиабилетов реализован посредством перехода на внешнюю платформу, поэтому отображение цен и бронирование осуществляются за пределами приложения.

Справочник городов, хотя и дополняется динамически через сервис подсказок, не охватывает все возможные населённые пункты. Эти ограничения не являются критическими и обусловлены главным образом отказом от коммерческих интерфейсов в пользу открытых данных.

Перспективы развития приложения связаны с расширением справочника направлений, добавлением возможности поиска обратного маршрута и многокомпонентных поездок, интеграцией дополнительных источников данных о транспорте и размещении, а также реализацией собственной базы данных для хранения истории поиска и пользовательских предпочтений.

Дополнительным направлением является локализация интерфейса на другие языки и адаптация приложения под требования конкретных образовательных организаций.

Таким образом, в ходе выполнения курсовой работы были решены все задачи, поставленные в задании. Изучена научная и техническая литература по вопросам проектирования веб-приложений, проведён анализ рынка цифровых туристических сервисов, рассмотрены доступные источники данных, разработана архитектура и реализовано работающее приложение Travel with Ease, выполнено его тестирование и сформулированы рекомендации по развитию.

Достигнута поставленная цель — теоретически обосновано и практически реализовано веб-приложение для поиска авиабилетов и бронирования отелей, адаптированное для русскоязычных пользователей и функционирующее без использования платных ключей доступа.

\phantomsection\addcontentsline{toc}{chapter}{Список литературы}

\chapter*{Список литературы}

\begin{enumerate}

  \item Андреева, Г. М. Психология общения : учебное пособие для вузов / Г. М. Андреева. — Москва : Аспект Пресс, 2012. — 368 с. — Текст : непосредственный.

  \item Бодалев, А. А. Психология общения : избранные психологические труды / А. А. Бодалев. — 4-е изд., перераб. и доп. — Москва : МПСИ ; Воронеж : МОДЭК, 2008. — 318 с. — Текст : непосредственный.

  \item Коджаспирова, Г. М. Педагогика в схемах и таблицах : учебное пособие / Г. М. Коджаспирова. — Москва : Проспект, 2016. — 248 с. — Текст : непосредственный.

  \item Леонтьев, А. А. Педагогическое общение / А. А. Леонтьев. — 2-е изд., перераб. и доп. — Москва : НИИ школьных технологий, 2009. — 368 с. — Текст : непосредственный.

  \item Петровская, Л. А. Компетентность в общении: социально-психологический тренинг / Л. А. Петровская. — Москва : Изд-во Моск. ун-та, 1989. — 216 с. — Текст : непосредственный.

  \item Сергеева, И. В. Информационные технологии в образовании : учебное пособие / И. В. Сергеева. — Москва : Академия, 2015. — 240 с. — Текст : непосредственный.

  \item Гаврилов, М. В. Веб-программирование : учебник для вузов / М. В. Гаврилов. — Москва : Юрайт, 2021. — 412 с. — Текст : непосредственный.

  \item Федоров, А. В. Проектирование информационных систем : учебное пособие / А. В. Федоров. — Санкт-Петербург : Питер, 2019. — 288 с. — Текст : непосредственный.

  \item OpenSky Network. API Documentation [Электронный ресурс]. — Режим доступа: \url{https://opensky-network.org/apidoc/} (дата обращения: 11.12.2025). — Текст : электронный.

  \item OpenStreetMap. About OpenStreetMap [Электронный ресурс]. — Режим доступа: \url{https://www.openstreetmap.org/about} (дата обращения: 11.12.2025). — Текст : электронный.

  \item Open-Meteo. Weather Forecast API Documentation [Электронный ресурс]. — Режим доступа: \url{https://open-meteo.com/en/docs} (дата обращения: 11.12.2025). — Текст : электронный.

  \item Aviasales. Поиск авиабилетов [Электронный ресурс]. — Режим доступа: \url{https://www.aviasales.ru} (дата обращения: 11.12.2025). — Текст : электронный.

  \item FastAPI. FastAPI Documentation [Электронный ресурс]. — Режим доступа: \url{https://fastapi.tiangolo.com} (дата обращения: 11.12.2025). — Текст : электронный.

\end{enumerate}

\phantomsection\addcontentsline{toc}{chapter}{Приложение}

\chapter*{Приложение}

\phantomsection\addcontentsline{toc}{section}{Приложение 1}

\section*{Приложение 1}

\begin{table}[htbp]\centering

\textbf{Фрагмент справочника кодов городов (IATA)}\par\vspace{4pt}

\begin{center}
\small
\begin{tabularx}{\textwidth}{|X|X|X|X|}
\hline
Город & Код & Город & Код \\
\hline
Москва & MOW & Санкт-Петербург & LED \\
\hline
Сочи & AER & Казань & KZN \\
\hline
Екатеринбург & SVX & Новосибирск & OVB \\
\hline
Владивосток & VVO & Краснодар & KRR \\
\hline
Калининград & KGD & Минеральные Воды & MRV \\
\hline
Дубай & DXB & Стамбул & IST \\
\hline
Бангкок & BKK & Пекин & BJS \\
\hline
\end{tabularx}
\end{center}

\end{table}

\phantomsection\addcontentsline{toc}{section}{Приложение 2}

\section*{Приложение 2}

Пример сформированной глубинной ссылки

Исходные данные:

\begin{enumerate}

  \item Город вылета: Москва (MOW)

  \item Город прибытия: Санкт-Петербург (LED)

  \item Дата поездки: 15.12.2026

  \item Количество пассажиров: 2

\end{enumerate}

Итоговая ссылка:

\url{https://www.aviasales.ru/search/MOW1512LED2}

Пояснение: строка «MOW1512LED2» содержит код города вылета (MOW), дату в формате ДДММ (1512), код города прибытия (LED) и количество пассажиров (2).

\phantomsection\addcontentsline{toc}{section}{Приложение 3}

\section*{Приложение 3}

\par\noindent\textbf{Пример ответа серверного интерфейса /api/flight\_link}\par\vspace{4pt}

\begin{Verbatim}[fontsize=\small,frame=single,xleftmargin=6pt]

{
  "from_city": "Москва",
  "to_city": "Санкт-Петербург",
  "from_code": "MOW",
  "to_code": "LED",
  "date": "2026-12-15",
  "adults": 2,
  "url": "https://www.aviasales.ru/search/MOW1512LED2",
  "note": "Полное расписание, время вылета и реальные цены — на Авиасейлс"
}

\end{Verbatim}

\phantomsection\addcontentsline{toc}{section}{Приложение 4}

\section*{Приложение 4}

\par\noindent\textbf{Фрагмент серверного модуля на платформе FastAPI (формирование ссылки на поиск рейсов)}\par\vspace{4pt}

\begin{Verbatim}[fontsize=\small,frame=single,xleftmargin=6pt]

@app.get("/api/flight_link")
def flight_link(
    from_city: str = Query(...),
    to_city: str = Query(...),
    date: str = Query(...),
    adults: int = Query(1),
):
    dep = city_iata(from_city)
    arr = city_iata(to_city)
    if not dep or not arr:
        return JSONResponse(
            {"error": "Город не найден"}, status_code=400
        )
    d, m = date.split("-")[2], date.split("-")[1]
    deep = f"{dep}{d}{m}{arr}{adults}"
    url = f"https://www.aviasales.ru/search/{deep}"
    return {"from_code": dep, "to_code": arr, "url": url}

\end{Verbatim}

\phantomsection\addcontentsline{toc}{section}{Приложение 5}

\section*{Приложение 5}

\par\noindent\textbf{Фрагмент клиентского модуля на JavaScript (обработка запроса поиска)}\par\vspace{4pt}

\begin{Verbatim}[fontsize=\small,frame=single,xleftmargin=6pt]

async function searchFlights() {
  const from = document.getElementById("f-from").value.trim();
  const to = document.getElementById("f-to").value.trim();
  const date = document.getElementById("f-date").value;
  const adults = document.getElementById("f-adults").value;
  const r = await fetch("/api/flight_link?from_city=" + encodeURIComponent(from) +
    "&to_city=" + encodeURIComponent(to) + "&date=" + date + "&adults=" + adults);
  const data = await r.json();
  if (data.error) { showError(data.error); return; }
  renderResult(data);
}

\end{Verbatim}

\end{document}
